\documentclass[a4paper, 12pt]{report}
\usepackage[hidelinks]{hyperref}
\usepackage[a4paper, margin=2cm]{geometry}
\usepackage[english, russian]{babel}
\usepackage{graphicx}
\usepackage{wrapfig}
\usepackage{amsmath}
\usepackage{amssymb}
\usepackage{xcolor}
\usepackage{minted}


\setlength{\parindent}{0em}
\setlength{\parskip}{0pt}

\begin{document}

\begin{titlepage}
    \centering
    {\large \textbf{ФЕДЕРАЛЬНОЕ ГОСУДАРСТВЕННОЕ АВТОНОМНОЕ ОБРАЗОВАТЕЛЬНОЕ УЧРЕЖДЕНИЕ ВЫСШЕГО ОБРАЗОВАНИЯ} \par}
    {\large «НАЦИОНАЛЬНЫЙ ИССЛЕДОВАТЕЛЬСКИЙ УНИВЕРСИТЕТ ИТМО» \par}
    \vspace{9cm}
    {\large {Отчёт по лабораторной работе} \par}
    {\large \textbf{«Типовая работа по математическому анализу №1»} \par}

    \vfill 
    \begin{flushright}
        {\large Выполнил работу: \par}
        {\large Косарев И. А. \par}
        {\large Академическая группа: \par}
        {\large № J3110 \par}
        \vspace{1cm}
        {\large Принято \par}
        {\large Табиева А. В. \par}
    \end{flushright}
    \vspace{1cm}
    {\large Санкт-Петербург \\ 2025}
\end{titlepage}

\tableofcontents
\newpage

\chapter*{Введение}
\addcontentsline{toc}{chapter}{Введение}
\setcounter{chapter}{1}
\setcounter{section}{0}

\section{Цели и задачи}
Цель: исследование интегрируемости функции $f(x) = e^{x}$ на промежутке $[0;1]$.

\subsection{Задачи}

\begin{itemize}
    \item Выполнить анализ функции: монотонность, ограниченность, интегрируемость
    \item Реализовать алгоритм взятия определенного интеграла
    \item Вычислить отклонение от истинного значения для каждого способа вычисления интеграла и сравнить
\end{itemize}

\chapter*{Аналитический этап}
\addcontentsline{toc}{chapter}{Аналитический этап}
\setcounter{chapter}{2}
\setcounter{section}{0}

\section{Верхняя и нижняя суммы Дарбу}
\textbf{Предварительный анализ} \par
Заметим, что $f^\prime(x) = (e^{x})^\prime = e^{x} > 0$ $ \forall x \in [0;1]$, что означает, что 
$f(x)$ возрастает на всем отрезке. Это значит, что $M_{i} = sup(x_i, x_{i+1})_{f} = f(x_{i+1})$, а
$m_i = inf(x_i, x_{i+1})_f = f(x_{i})$. Тогда верхняя и нижняя суммы Дарбу соответственно будут иметь 
следующий вид (берем конгруэнтное разбиение отрезка на n частей):
\[
U(f, t) = \sum_{i=0}^{n-1}\Delta{x_i}M_i = \sum_{i=0}^{n-1}\frac{1}{n}e^{x_{i+1}} = \sum_{i=0}^{n}\frac{1}{n}e^{\frac{i+1}{n}}
\]
\[
L(f, t) = \sum_{i=0}^{n-1}\Delta{x_i}m_i = \sum_{i=0}^{n-1}\frac{1}{n}e^{x_i} = \sum_{i=0}^{n}\frac{1}{n}e^{\frac{i}{n}}
\]

\textbf{Вычисление сумм} \par
Вычислим обе суммы по формуле суммы геометрической прогрессии:
\[
U(f, t) = \frac{1}{n}\sum_{i=0}^{n-1}e^{\frac{1}{n}}(e^{\frac{1}{n}})^{i} = \frac{1}{n}\frac{e^{\frac{1}{n}}(1 - e^{\frac{1}{n}n})}{1 - e^{\frac{1}{n}}}
\]
\[
L(f, t) = \frac{1}{n}\sum_{i=0}^{n-1}(e^{\frac{1}{n}})^{i} = \frac{1}{n}\frac{(1 - e^{\frac{1}{n}n})}{1 - e^{\frac{1}{n}}}
\]

\section{Критерий Римана интегрируемости функции}
Определение критерия:
\[
f(x) \in R([a,b]) \Leftrightarrow \forall \epsilon > 0 \; \exists \tau: U(f, \tau) - L(f, \tau) < \epsilon
\]

Рассмотрим: 
\[
U(f, \tau) - L(f, \tau) = (e^{\frac{1}{n}} - 1)\frac{1}{n}\frac{(1 - e)}{1 - e^{\frac{1}{n}}} = \frac{1}{n}(e - 1) < \epsilon
\]
Возьмем $n = [\frac{e-1}{\epsilon}] + 1$. При данном $n$ неравенство выполняется, значит, $f(x) = e^x \in R([0,1])$.

\section{Интегралы Дарбу}
Рассмотрим (для устремления ранга к нулю берем конгруэнтное разбиение):
\[
U(f) = \lim_{\lambda(\tau) \rightarrow 0}U(f, \tau) = \lim_{n \rightarrow \infty} \frac{\frac{1}{n}e^{\frac{1}{n}}(1 - e)}{1 - e^{\frac{1}{n}}} = [\frac{0}{0}] = \lim_{n \rightarrow \infty} (e - 1)\frac{\frac{1}{n^2}e^{\frac{1}{n}} + \frac{1}{n^3}e^{\frac{1}{n}}}{e^{\frac{1}{n}}\frac{1}{n^2}} = \lim_{n \rightarrow \infty} (e - 1)(1 + \frac{1}{n}) = e - 1
\]

\[
L(f) = \lim_{\lambda(\tau) \rightarrow 0}L(f, \tau) = \lim_{n \rightarrow \infty} \frac{\frac{1}{n}(1 - e)}{1 - e^{\frac{1}{n}}} = [\frac{0}{0}] = \lim_{n \rightarrow \infty} \frac{(e - 1)\frac{1}{n^2}}{e^{\frac{1}{n}}\frac{1}{n^2}} = \lim_{n \rightarrow \infty} (e - 1) e^{-\frac{1}{n}} = e - 1
\]
Как можно заметить, интегралы совпадают, что также является подтверждением $f(x) = e^x \in \, R([0, 1])$. 

\section{Доп. достаточное условие интегрируемости}
Для дополнительной проверки интегрируемости функции $f(x)$ можно воспользоваться критерием Лебега (или же следствием из него). $f(x) = e^x$, как говорилось ранее, монотонна на всем промежутке, из чего следует, что
\[
\forall \; x \in [0,1] \; f(0) = 1 \leq e^x \leq e = f(1)
\]
Значит, функция ограничена и монотонна $\rightarrow$ по следствию из критерия Лебега функция интегрируема.

\section{Сравнение значения с формулой Ньютона-Лейбница}
Вычислим неопределенный интеграл от $f(x)$:
\[
\int f(x)\,dx = e^x + C
\]
Вычислим определенный интеграл по формуле Ньютона-Лейбница:
\[
\int_{0}^{1} f(x)\,dx = e^x \big|_{0}^{1} = e - 1
\]
Как можно заметить, значение совпало с предыдущими вычислениями.

\chapter*{Практический этап}
\addcontentsline{toc}{chapter}{Практический этап}
\setcounter{chapter}{3}
\setcounter{section}{0}

\section{Реализация}
Были написаны алгоритмы вычисления определенного интеграла:

\begin{itemize}
    \item Методом прямоугольников (левые, правые, середина)
    \item Методом трапеций
    \item Методом Симпсона
\end{itemize}

Также были реализованы вспомогательные функции для конгруэнтного и рандомизированного резбиения отрезка и для визуализации сравнения точности всех методов (принцип работы описан в теле функций).

\section{Код реализации}
\inputminted[frame=single,fontsize=\small]{python3}{./utils.py}

\section{Анализ и построение графиков}
Графики и их построение приведены в файле Jupyter Notebook, сохраненном в \textcolor{red}{\href{https://github.com/quant1on/matan_lab_2}{репозитории лабораторной}}. Ниже приведен график сравнения ошибки разных методов интегрирования:

\begin{figure}[h]
    \centering
    \includegraphics[width=0.45\textwidth]{image.png}
    \caption{\small{Сравнение MAE (логарифмированной) в зависимости от точности интегрирования (кол-ва подотрезков)}}
\end{figure} \par

За основу вычисления ошибки был взят принцип MAE. Как видно на графике, самым точным методом оказался метод Симпсона. Можно заметить, что при увеличении кол-ва подотрезков интегрирования ошибка для всех методов уменьшается, что говорит о прямой корреляции точности и кол-ва подотрезков (ранга разбиения).

\section*{Выводы}
\addcontentsline{toc}{chapter}{Выводы}
\setcounter{chapter}{4}
\setcounter{section}{0}
В процессе выполнения лабораторной работы был выполнен анализ функции $f(x) = e^x$, проверка ее свойств. Также был написан простейший алгоритм интегрирования разными методами. Программа позволяет выбрать метод взятия границ подотрезков и их количество, что позволяет регулировать точность вычисления. \par
Также было визуализировано сравнение точности каждого из методов и проведен анализ данного графика. \par
В процессе выполнения лабораторной работы я смог получше разобраться с вычислением интегральных сумм и смог на практике применить свои теоретические знания.

\end{document}